\subsection{Implications}
\label{subec:implications}

The discussion above leads to three broad sets of theoretical and substantive implications.  First, consistent with the discussion in~\ref{subsec:relform}, we should expect that those individuals who have a larger propensity to engage in culture-talk will also have an easier time during the relationship-formation process.  This will then be a plausible mechanism to account for the existence stubborn status-linked inequalities social capital.  For instance, it is well-known that education has been \citep{fischer82, marsden87} and continues to be \citep{lizardo06a, diprete_etal11} one of the best predictors of network size in the United States (and other national settings) today.  If the higher-educated are also more likely to engage in culture-talk than the less-educated (within-levels of consumption), and thus are more likely to routinely use this as an interactional resource to form new relationships, then this would serve as a plausible micro-mechanism accounting for this effect.  The same can be said for those socio-demographic groups known to have smaller networks than average (e.g. the elderly, and some ethnic and racial minorities).  If culture-talk propensities decline with age or are different across racial and ethnic groups, then we should also expect rates of relationship formation to track these differences.

Second, in line with the formulation in~\ref{subsec:solidarity}, we should expect that persons who belong to socio-demographic groups that exhibit a larger propensity to engage in culture-talk will also tend to be embedded in status segregated networks.   Social ties that are partially constituted via culture-talk retain this property because they are dynamically sustained via conversational meanings and contents access to which is also status-segregated \citep{bourdieu84}.   In this respect, differential rates of culture-talk across socio-demographic groups emerge as a plausible mechanism that accounts for the existence of status-segregated social networks even in the absence of any strictly conscious attempt on the part of individuals to select status-homogamous contacts.   It is well-established that the social-networks of Americans are characterized by high-rates of status-endogamy specially premised on race and education \citep{fischer82, marsden87, mcpherson_etal06, diprete_etal11}, and that marital homogamy premised on education and income has increased at the expense of marital endogamy premised on other axes of difference \citep{kalmijn91, kalmijn94, schwartz_mare05, schwartz10}.  If persons who belong to high-education or high-income groups are also more likely to engage in culture-talk than the less-educated (within-levels of consumption), then they should also be expected to be more likely to routinely use this as an interactional resource to engage in status-segregated conversational rituals.

Finally, following the discussion in~\ref{subsec:linkages}, we should expect that those individuals who have a more pronounced propensity to engage in culture-talk will also tend to be more likely to bridge structural and ``cultural holes'' \citep{pachucki_breiger10, schultz_breiger10, lizardo06a}.   Status-linked inequalities in culture-talk propensities can then be proposed as a mechanism that explains the often noted fact that high-status individuals are more likely to be connected to a more diverse pool of others than low status individuals \citep{marsden87, erickson96, erickson08, lin_dumin86, lin99, goldstein_warren00, diprete_etal11}.   If persons who belong to the higher educated and those of high-income groups are also more likely to engage in culture-talk than the less-educated (within-levels of consumption) and those of low income, and are able to (in some settings) use this cultural expertise to make sporadic connections \emph{across} group and status boundaries, then differential propensities to engage in culture talk emerge as a natural micro-level explanation for this phenomenon. 
 